% Общий глоссарий всех пяти документов. Один термин — одна строка \term{Ключ}{Определение}, по алфавиту.
% Документ печатает нужные: \printterms{ЕК, СП} или все: \printallterms.
% Термин меняется только здесь, тогда он одинаковый во всех PDF. Дополняется по промпту 5.7 из README.md.
\term{Базовая версия}{утверждённая версия единицы конфигурации; промежуточные версии базовыми не считаются (слайды 75--76).}
\term{ЕК}{единица конфигурации: элемент под управлением конфигурацией, пара «уникальное имя, версия». Планы и стандарты тоже ЕК.}
\term{ЗИ}{запрос на изменение: описывает планируемое изменение одной ЕК и порождается из утверждённого СП (слайды 73--74).}
\term{КТ-178}{русский перевод авиационного стандарта DO-178 о разработке бортового ПО; на нём построен курс.}
\term{КТК}{контрольная точка качества: проверка, которую релиз проходит перед продвижением дальше. Сокращение «КТ» занято стандартом КТ-178.}
\term{СП}{сообщение о проблеме: регистрирует одно несоответствие в одной или нескольких ЕК (слайды 64--66).}
